\subsection{Quantile and Shape Displays}
\label{subsec:quantile-and-shape-displays}
% Source integration: attached screenshots from Chapter 10, "Summarizing Data."

Distributional shape can be read from cumulative displays, density displays, or quantile comparisons.  These views emphasize different questions: cumulative proportion, local concentration, and relative change between two distributions.

\begin{definition} \label{def:qq-plot} \hlt{Quantile--Quantile Plot}\\
Let $F$ and $G$ be continuous strictly increasing CDFs with quantile functions as in Definition \ref{def:quantile-and-equality-in-distribution}.  A \hlt{Q--Q plot} compares the two distributions by plotting
\[
\bigl(F^{-1}(p),G^{-1}(p)\bigr),\qquad 0<p<1.
\]
For two equal-sized batches with order statistics $X_{(1)}\leq\cdots\leq X_{(n)}$ and $Y_{(1)}\leq\cdots\leq Y_{(n)}$, the empirical Q--Q plot uses the points $(X_{(i)},Y_{(i)})$.
\end{definition}

\begin{proposition} \label{prop:qq-location-scale-effects} \hlt{Location and Scale Effects in Q--Q Plots}\\
Let $x_p=F^{-1}(p)$ and $y_p=G^{-1}(p)$.
\begin{enumerate}[label=(\roman*)]
\setlength{\itemsep}{0pt}
\item If $y_p=x_p+h$ for all $p$, then $G(y)=F(y-h)$ and the Q--Q plot is a line of slope $1$ and intercept $h$.
\item If $y_p=cx_p$ for all $p$ with $c>0$, then $G(y)=F(y/c)$ and the Q--Q plot is a line through the origin with slope $c$.
\end{enumerate}
\begin{proof}
For (i), $p=G(y_p)=F(x_p)=F(y_p-h)$ for every $p$, so the CDFs agree by reparameterizing $y=y_p$.  The proof of (ii) is the same with $x_p=y_p/c$.
\end{proof}
\end{proposition}

\begin{figure}[H]
  \centering
  % Source: supplied Chapter 10 screenshots on Q--Q plots and simple treatment effects.
% Schematic redraw, not a reproduction of the scanned figures.
\begin{tikzpicture}[
  x=0.72cm,
  y=0.72cm,
  >=Stealth,
  line cap=round,
  line join=round,
  every node/.style={font=\scriptsize},
  axis/.style={->, gray!70, line width=0.45pt},
  add/.style={blue!70!black, line width=0.95pt},
  mult/.style={green!45!black, line width=0.95pt},
  ref/.style={gray!55, line width=0.55pt},
  point/.style={black, fill=black, inner sep=0pt, minimum size=2pt, circle}
]
  \begin{scope}
    \draw[axis] (0,0) -- (4.6,0) node[black!78,right] {$x_p$};
    \draw[axis] (0,0) -- (0,4.6) node[black!78,above] {$y_p$};
    \draw[ref] (0,0) -- (4.1,4.1);
    \draw[add] (0,0.72) -- (3.72,4.44);
    \foreach \x in {.55,1.25,1.95,2.65,3.35}
      \node[point] at (\x,{\x+.72}) {};
  \end{scope}

  \begin{scope}[xshift=5.6cm]
    \draw[axis] (0,0) -- (4.6,0) node[black!78,right] {$x_p$};
    \draw[axis] (0,0) -- (0,4.6) node[black!78,above] {$y_p$};
    \draw[ref] (0,0) -- (4.1,4.1);
    \draw[mult] (0,0) -- (2.96,4.44);
    \foreach \x in {.45,.95,1.45,1.95,2.45}
      \node[point] at (\x,{1.5*\x}) {};
  \end{scope}
\end{tikzpicture}

  \caption{Straight Q--Q plots diagnose simple additive and multiplicative effects.}
  \label{fig:eda-qq-effects}
\end{figure}

\begin{example} \label{ex:normal-and-exponential-qq-lines} \hlt{Normal and Exponential Q--Q Lines}\\
If $F$ is $N(0,1)$ and $G$ is $N(1,1)$, then $y_p=x_p+1$.  If $G$ is instead $N(1,4)$, then $y_p=1+2x_p$.  Thus the first Q--Q plot has slope $1$ and intercept $1$, while the second has slope $2$ and intercept $1$.\\
If $F$ is exponential with rate $\lambda_1$ and $G$ is exponential with rate $\lambda_2$, then
\[
F^{-1}(p)=-\frac{\log(1-p)}{\lambda_1},
\qquad
G^{-1}(p)=-\frac{\log(1-p)}{\lambda_2},
\]
so the Q--Q plot is the line $y=(\lambda_1/\lambda_2)x$.
\end{example}

\begin{method} \label{meth:interpolated-empirical-quantile} \hlt{Interpolated Empirical Quantiles}\\
For order statistics $X_{(1)},\ldots,X_{(n)}$, assign $X_{(k)}$ to probability $k/(n+1)$.  If
\[
\frac{k}{n+1}\leq p\leq\frac{k+1}{n+1},
\]
one linear interpolation convention is
\[
\widehat{x}_p
=(1-\lambda)X_{(k)}+\lambda X_{(k+1)},
\qquad
\lambda=(n+1)p-k.
\]
This lets Q--Q plots compare unequal batch sizes by fixing a common grid of probabilities and interpolating the larger or smaller sample as needed.
\end{method}

\begin{proposition} \label{prop:sample-quantile-variance-delta-method} \hlt{Approximate Variance of a Sample Quantile}\\
Let $Y_1,\ldots,Y_n$ be iid from a continuous CDF $F$ with density $f$, and let $x_p=F^{-1}(p)$ with $f(x_p)>0$.  For an order statistic whose rank is near $p(n+1)$,
\[
\operatorname{Var}(\widehat{x}_p)
\approx
\frac{p(1-p)}{n f(x_p)^2}.
\]
\begin{proof}
By Theorem \ref{thm:kth-order-statistic-density}, the transformed order statistic $F(Y_{(k)})$ behaves like a uniform order statistic with mean approximately $p$ and variance approximately $p(1-p)/n$.  Applying the derivative of $F^{-1}$, namely $1/f(F^{-1}(p))$, gives the displayed delta-method approximation.
\end{proof}
\end{proposition}

\begin{definition} \label{def:histogram-density-and-stem-leaf-display} \hlt{Histogram, Density Display, and Stem-and-Leaf Plot}\\
\begin{enumerate}[label=(\roman*)]
\setlength{\itemsep}{0pt}
\item A \hlt{histogram} divides the range into bins and plots counts or proportions in each bin.  With unequal bin widths, density height should be proportional to ``proportion divided by bin width'' so that area represents probability.
\item A \hlt{density display} estimates local concentration by a smooth curve, often using the kernel estimator from Definition \ref{def:kernel-density-estimator}.
\item A \hlt{stem-and-leaf plot} writes each observation as a common leading part, the stem, and a final digit or block, the leaf.  It displays shape while retaining much of the original numerical information.
\end{enumerate}
\end{definition}

\begin{figure}[H]
  \centering
  % Source: supplied Chapter 10 screenshots on histograms and kernel density estimates.
% Schematic redraw, not a reproduction of the scanned figures.
\begin{tikzpicture}[
  x=0.62cm,
  y=0.42cm,
  >=Stealth,
  line cap=round,
  line join=round,
  every node/.style={font=\scriptsize},
  axis/.style={->, gray!70, line width=0.45pt},
  bar/.style={draw=blue!70!black, fill=blue!10, line width=0.45pt},
  rough/.style={blue!70!black, line width=0.95pt},
  wideband/.style={green!45!black, line width=0.95pt}
]
  \begin{scope}
    \draw[axis] (0,0) -- (8.2,0) node[black!78,right] {$x$};
    \draw[axis] (0,0) -- (0,5.5) node[black!78,above] {count};
    \foreach \lo/\hi/\h in {
      .45/1.05/1, 1.05/1.65/2.4, 1.65/2.25/2,
      2.25/2.85/4.7, 2.85/3.45/3.8, 3.45/4.05/4.7,
      4.05/4.65/2.4, 4.65/5.25/3, 5.25/5.85/1.1,
      5.85/6.45/2
    }{
      \draw[bar] (\lo,0) rectangle (\hi,\h);
    }
  \end{scope}

  \begin{scope}[xshift=6.35cm]
    \draw[axis] (0,0) -- (8.2,0) node[black!78,right] {$x$};
    \draw[axis] (0,0) -- (0,5.5) node[black!78,above] {$\widehat f_h$};
    \draw[rough]
      plot[smooth] coordinates {
        (.45,.25) (.85,.75) (1.2,.45) (1.55,1.85) (1.9,1.1)
        (2.25,3.1) (2.6,2.0) (2.95,4.45) (3.3,3.05)
        (3.65,4.55) (4.0,3.25) (4.45,3.85) (4.9,2.35)
        (5.35,2.8) (5.85,1.45) (6.35,.7) (6.95,.3)
      };
    \draw[wideband, domain=.35:7.3, samples=120]
      plot (\x,{4.05*exp(-0.18*(\x-3.45)*(\x-3.45))});
  \end{scope}
\end{tikzpicture}

  \caption{Bin width and bandwidth control the roughness of shape displays.}
  \label{fig:eda-histogram-kde-bandwidth}
\end{figure}

\begin{remark} \label{rem:shape-displays-lose-different-information} \hlt{What Shape Displays Preserve}\\
An empirical CDF preserves all order information and supports direct quantile reading.  A histogram sacrifices exact data values for a compact view of shape.  A kernel estimate smooths the histogram idea, but the bandwidth controls how much local structure is retained.  A stem-and-leaf plot is useful when one wants both a rough distributional shape and recoverable numerical values.
\end{remark}
